\section{Test Case 6: Three-D simulation of a quasi-Newtonian fluid}

The input in the \textit{input-6} sub-directory activates the
Kelvin--Voigt constitutive model (\texttt{kvfluid yes},
Eq.~\ref{eq:stress-ode-KV}) in place of the Maxwell model, and provides the
reference case for the quasi-Newtonian/Kelvin--Voigt integration path.  It
keeps the geometry of Test Case 3 (Sec.~\ref{sec:Test-Case-3:}) with the
Heun integrator (\texttt{integrator 2}), a viscosity of 10 P, a negligible
elastic modulus ($10^{-2}$ dyn/cm$^{2}$), twice the charge density
(88000 statC/cm$^{3}$), and the multiple timestep algorithm for Coulomb
forces (Sec.~\ref{sec:Multiple-Timestep-Algorithm}).  The OpenACC build runs
this case with the CPU build's code, since its device step does not cover
multiple-step Coulomb sums yet.

Over $10^{6}$ steps (0.01 s) the multiple timestep algorithm agrees with the
direct summation within $9\times10^{-8}$ in the axial position and
$4.4\times10^{-6}$ in the off-axis distance, with the same 146 insertions
and 35 removals (one of them a step later), and two MPI ranks agree with the
serial run within $6\times10^{-6}$.  Versions up to 1.22 departed from the
direct summation by up to 25\% in the off-axis distance within the same
window, with one more insertion, because of the defects corrected in
version 2.0 (Sec.~\ref{sec:Multiple-Timestep-Algorithm}); their per-rank
arrays were also one element short, and once the jet filled its arrays a
GFortran build stopped near step 668,500 on an out-of-bounds write, which
other builds made silently.
